\section{Robust testing: complete constructions and proofs}
\label{sec:annex-p15-robust-hypothesis-testing}

\begin{annexitem}{annex:p15.robust.overlap}{thm:p15-robust-overlap}{Measure domination, variation and intersecting classes}
\begin{proof}
First record the variation identity used throughout. For probabilities
\(P,R\), take \(\nu=P+R\) and densities \(p,r\). Their difference
has integral zero, so
\[
 \int(p-r)_+d\nu=\int(r-p)_+d\nu
 =\tfrac12\int|p-r|d\nu.
\]
Every event difference is bounded above by the first positive-part
integral, attained on \(\{p\geq r\}\); the negative difference has the
same maximum. Thus this quantity is \(d_{\rm TV}(P,R)\), and
\(\|P-R\|_{\rm var}=\int|p-r|d\nu\).
The sup-event formula proves the triangle inequality. It also shows
that the definition is independent of a dominating measure.

A member \(R\) of a contamination class satisfies \(R\geq a_iP_i\)
as measures. Conversely any probability with this domination has
\(R-a_iP_i\) a positive measure of mass \(\varepsilon_i\), giving
the required contaminant when \(\varepsilon_i>0\). If \(\varepsilon_i=0\),
the positive difference has mass zero and \(R=P_i\).
If \(R\) belongs to both classes, use the common finite dominating
measure \(R+P_0+P_1\). Its density is at least both \(a_ip_i\), so its
unit integral is at least \(m\). Changing from this dominating measure
to \(\mu=P_0+P_1\) gives the same maximum integral: both nominal
densities gain the common factor \(d\mu/d(R+P_0+P_1)\).
Thus \(m\leq1\) is necessary.
If \(m\leq1\), the measure
\[
 R(dx)=\max(a_0p_0(x),a_1p_1(x))\mu(dx)+(1-m)P_0(dx)
\]
is a probability dominating both \(a_iP_i\), proving sufficiency.
For equal contamination,
\(\int\max(p_0,p_1)d\mu=1+d\), by the variation identity.
Hence \(m=(1-e)(1+d)\), giving the stated separation threshold.

For variation balls, a common member requires
\(\|P_1-P_0\|_{\rm var}\leq2\varepsilon\), so \(\varepsilon\geq d\).
Conversely \((P_0+P_1)/2\) is at variation distance \(d\) from either
nominal law and belongs to both balls when \(\varepsilon\geq d\).

If \(R\) belongs to both classes, then \(R^{\otimes n}\) is allowed
under either hypothesis. A level-\(\alpha\) test has power at that
alternative at most \(\alpha\); the constant test attains exactly that
worst-case power. For any test put \(z=E_{R^{\otimes n}}\varphi\).
Its maximum risk is at least
\(\max\{L_0z,L_1(1-z)\}\), minimized over \(0\leq z\leq1\) when
the two terms agree. This gives \(z=L_1/(L_0+L_1)\) and the stated
value, attained uniformly by the constant test.
Similarly its weighted worst-case risk is at least
\(w_0z+w_1(1-z)\geq\min(w_0,w_1)\).
An always-accept or always-reject decision attains the smaller value.
\end{proof}
\end{annexitem}

\begin{annexitem}{annex:p15.huber.contamination.pair}{thm:p15-huber-contamination-pair}{Existence, uniqueness and the source's normalized densities}
\begin{proof}
The elementary identity
\[
 F(u)=1+\int(up_0-p_1)_+d\mu
\]
shows \(F(0)=1\), \(F(u)\geq u\), and
\(0\leq F(u+\Delta)-F(u)\leq\Delta\) for \(\Delta\geq0\).
Thus \(F\) is continuous, nondecreasing and unbounded.
If \(u_0\) is the essential infimum of \(r\) under \(P_0\), its
positive-part formula says \(F(u)=1\) exactly for \(0\leq u\leq u_0\).
Once \(F(u)>1\), the set \(\{up_0>p_1\}\) has positive \(P_0\) mass,
and
\[
 F(u+\Delta)-F(u)\geq
        \Delta P_0\{up_0>p_1\}>0\quad(\Delta>0).
\]
This proves existence and uniqueness at every target greater than one.
The zero-contamination target one is satisfied by the specified choice
\(u=0\). Interchanging indices proves all assertions for \(G\).
These are the continuous positive-part arguments of the source's
Lemma 1, including its possibly nonunique zero-contamination case.

The equations make \(\int q_i d\mu=1\).
Also \(q_i\geq a_ip_i\), so the measure-domination argument in the
preceding proof gives \(Q_i\in\mathcal C_i\).
Their excess measures have densities
\[
 q_0-a_0p_0=a_0(vp_1-p_0)_+,\qquad
 q_1-a_1p_1=a_1(up_0-p_1)_+,
\]
of masses \(\varepsilon_0,\varepsilon_1\), respectively.
They place their mass in the upper and lower clipping regions.

Since \(b=a_1/a_0\),
\[
 a_1F(1/b)=m>1,\qquad a_0G(b)=m>1.
\]
The normalized targets have value one after these multipliers.
Monotonicity and the specified zero-target choices imply \(u<1/b\)
and \(v<b\). Thus \(u<1/v\) whenever \(v>0\), and the same extended
inequality holds when \(v=0\).
For \(p_0,p_1>0\),
\[
 \frac{q_1}{q_0}
 =b\,\frac{\max(r,u)}{\max(1,vr)}
 =
 \begin{cases}
 bu,&r\leq u,\\
 br,&u<r<1/v,\\
 b/v,&r\geq1/v
 \end{cases}
\]
when \(v>0\). If \(v=0\), the last branch is absent.
The same calculation directly from the densities covers \(p_1=0<p_0\)
and \(p_0=0<p_1\), giving zero or infinity only in the stated
zero-contamination cases.
For example \(v>0,p_0=0<p_1\) gives \(q_1/q_0=b/v\).
On the common density-zero set both \(Q_i\) vanish, so the clipped
assignment does not change their laws but does define the test under
singular contaminations.

For \(0<c'<c''<\infty\), integrating the two branches for each density
gives the source's original equations explicitly:
\[
 a_0\{P_0(r<c'')+(c'')^{-1}P_1(r\geq c'')\}=1,\qquad
 a_1\{P_1(r>c')+c'P_0(r\leq c')\}=1.
\]
At zero contamination, the max-density formulas already proved give
the unambiguous endpoint interpretation without a product \(0\cdot\infty\).
\end{proof}
\end{annexitem}

\begin{annexitem}{annex:p15.huber.tv.pair}{thm:p15-huber-tv-pair}{The second neighborhood's mass transfer}
\begin{proof}
Let \(s=p_0+p_1\) and \(g(k)=\int(ks-p_1)_+d\mu\) on \([0,1]\).
We have \(g(0)=0\), \(g(1)=\int p_0d\mu=1\), and
\[
 0\leq g(k+\Delta)-g(k)\leq
             \Delta\int s\,d\mu=2\Delta .
\]
Thus \(g\) is continuous. Wherever \(g(k)>0\), the set
\(\{ks>p_1\}\) has positive \(s\,d\mu\) measure; increasing \(k\)
by \(\Delta>0\) increases \(g\) by at least
\(\Delta\int_{\{ks>p_1\}}s\,d\mu>0\).
Hence each target in \((0,1)\) has a unique solution.
The decreasing upper-transfer function
\(\int(p_1-ks)_+d\mu\) has the analogous property.
This proves the required form of Huber's Lemma 3.
Its initial flat region is exactly
\(0\leq k\leq\operatorname{ess\,inf}_{s\,d\mu}(p_1/s)\);
the ratio is considered only where \(s>0\).

At \(k=1/2\), both transfer integrals equal
\(\tfrac12\int(p_0-p_1)_+d\mu=d/2\).
Since \(0<\varepsilon/2<d/2\), the two solutions satisfy
\(0<k_-<1/2<k_+<1\).
The difference \(q_1-p_1\) is positive only in the lower region
and negative only in the upper region; its positive and negative
integrals both equal \(\varepsilon/2\).
Thus \(q_1\) integrates to one. Since \(q_0=s-q_1\) and \(\int s=2\),
it also integrates to one. Both are nonnegative.
The variation identity gives
\(\|Q_1-P_1\|_{\rm var}=\varepsilon\);
the equality \(q_0-p_0=-(q_1-p_1)\) gives the same value for index zero.

On the lower region \(q_1=k_-s,q_0=(1-k_-)s\), so their ratio is
\(c'=k_-/(1-k_-)\). The upper region gives \(c''\); the middle
region leaves the nominal densities unchanged.
These identities reproduce the three branches of the source's
equation (6), and the two transfer integrals are its equation (7).
At radius zero take the nominal pair directly; at or beyond the
overlap radius the preceding overlap theorem applies.
\end{proof}
\end{annexitem}

\begin{annexitem}{annex:p15.huber.stochastic.order}{thm:p15-huber-stochastic-order}{The full neighborhood order, including external support}
\begin{proof}
Write \(l,u\) for the lower and upper score endpoints.
Every point of the sample space, including a nominal null point, has
score between these endpoints. Thus for \(t\leq l\) the event
\(E_t=\{\ell<t\}\) is empty, and for \(t>u\) it is the whole space,
when the latter case is possible.
For \(l<t\leq u\), the contamination construction adds no mass to
\(Q_0\) on \(E_t\), while all the excess mass of \(Q_1\) is on \(E_t\).
Consequently
\[
 Q_0(E_t)=a_0P_0(E_t),\qquad Q_1(E_t)=a_1P_1(E_t)+\varepsilon_1.
\]
For arbitrary \(R_i=a_iP_i+\varepsilon_iH_i\), including singular \(H_i\),
\[
 R_0(E_t)\geq a_0P_0(E_t),\qquad
 R_1(E_t)\leq a_1P_1(E_t)+\varepsilon_1.
\]
This proves the two outer inequalities. Zero excess masses cause no
exception to this reasoning.

For variation neighborhoods, the lower region removes mass
\(\delta=\varepsilon/2\) from \(P_0\), while the upper region adds
that mass; the roles reverse for \(P_1\).
The event \(E_t\) includes the former entire region and excludes the
latter. Thus
\[
 Q_0(E_t)=P_0(E_t)-\delta,\qquad
 Q_1(E_t)=P_1(E_t)+\delta.
\]
The sup-event definition of variation gives the same two outer
comparisons for every \(R_i\) in its ball.

For the middle comparison use the density likelihood ordering.
If \(0<t\leq1\), then on \(E_t\) we have \(q_1\leq tq_0\), so
\(Q_1(E_t)\leq tQ_0(E_t)\leq Q_0(E_t)\).
If \(t>1\), on \(E_t^c\) we have \(q_1\geq tq_0\), including the
possible \(Q_1\)-only support. Hence
\(Q_1(E_t)\leq1-tQ_0(E_t^c)\leq Q_0(E_t)\).
Thresholds \(t\leq0\) give empty events.
Decreasing strict upper thresholds to a fixed threshold gives the
weak-threshold distribution inequalities by continuity of probability.
\end{proof}
\end{annexitem}

\begin{annexitem}{annex:p15.robust.product.order}{thm:p15-robust-product-order}{Quantile construction and the monotone rejection function}
\begin{proof}
Map the extended score to \(Z=\ell/(1+\ell)\in[0,1]\), assigning
the value one at infinity. The strict-threshold inequalities imply
the usual CDF inequalities for \(Z\), by the limits proved above.
For any distribution function \(F\) on \([0,1]\), and \(0<u<1\), define
\[
 q_F(u)=\inf\{z\in[0,1]:F(z)\geq u\}.
\]
The set is nonempty since \(F(1)=1\). Monotonicity and right
continuity show \(F(q_F(u))\geq u\), and therefore
\[
 q_F(u)\leq z\quad\Longleftrightarrow\quad u\leq F(z).
\]
The inverse is nondecreasing and measurable. For a uniform \(U\) on
\((0,1)\), the last equivalence gives
\(\Pr\{q_F(U)\leq z\}=F(z)\), proving its law.
If \(F\geq G\), their defining sets give \(q_F\leq q_G\).
Apply all four inverse functions to the same \(U\) to obtain the
ordered coupling; use independent uniforms for the \(n\) coordinates.
Each score vector then has precisely its independent product law.
This proof takes place entirely in the real score space.

When both contamination fractions are positive the scores lie in a
finite positive interval. When exactly one fraction vanishes, they
lie either in \([l,\infty]\) with \(l>0\), or in \([0,u]\) with
finite \(u\). The positive-radius variation case again has a finite
positive interval. Products are therefore defined and nondecreasing
in each coordinate in all these cases.
The function
\(z\mapsto{\bf1}_{\{z>c\}}+\kappa{\bf1}_{\{z=c\}}\) is nondecreasing
for \(0\leq\kappa\leq1\).
Pointwise comparison in the coupling, followed by expectation,
proves
\[
 E_{R_0^{\otimes n}}\varphi\leq E_{Q_0^{\otimes n}}\varphi,\qquad
 E_{R_1^{\otimes n}}\varphi\geq E_{Q_1^{\otimes n}}\varphi .
\]
The inequalities become attained suprema and infima because each
\(Q_i\) belongs to its neighborhood.
This proof also applies to the monotone test \({\bf1}_{\{T_n=\infty\}}\)
by increasing finite thresholds.
\end{proof}
\end{annexitem}

\begin{annexitem}{annex:p15.huber.minimax.tests}{thm:p15-huber-minimax-tests}{Neyman--Pearson, weighted risks and an equalizer}
\begin{proof}
Let \(g_i(x_1,\ldots,x_n)=\prod_jq_i(x_j)\), densities relative to
\(\mu^{\otimes n}\). On the support of these laws the likelihood
ordering is exactly that of \(T_n\). Where both product densities
vanish, the score-based test remains defined but does not affect
either simple-law risk.
For positive uncertainty in at least one class, the product conventions
in the preceding proof prevent the indeterminate expression
\(0\cdot\infty\). With zero uncertainty in both classes there is
no contamination to control outside the two simple laws; use their
product-density comparison directly.

For \(0<\alpha<1\), apply
\cref{thm:p15-likelihood-threshold,thm:p15-neyman-pearson}
to \(g_0,g_1\). Choose the constant randomization on the threshold
set, giving \(E_{Q_0^{\otimes n}}\varphi=\alpha\).
The product-order theorem says that its robust size is exactly
\(\alpha\) and its worst power is its power under \(Q_1^{\otimes n}\).
Any robustly feasible competitor \(\psi\) has simple-law size at most
\(\alpha\), so Neyman--Pearson gives
\[
 \inf_R E_{R^{\otimes n}}\psi
 \leq E_{Q_1^{\otimes n}}\psi
 \leq E_{Q_1^{\otimes n}}\varphi
 =\inf_R E_{R^{\otimes n}}\varphi ,
\]
where the infima are over the alternative class.
At level zero the simple optimal test rejects on \(g_0=0<g_1\),
which is \(T_n=\infty\) on the pair's support. Product order gives
its robust size zero and the worst-power comparison. At level one,
the constant-one test has power one under every law.

For positive \(w_0,w_1\), the simple weighted loss of any test is
\[
 \int\{w_0g_0\varphi+w_1g_1(1-\varphi)\}\,d\mu^{\otimes n}.
\]
The integrand is minimized pointwise by rejecting where
\(g_1/g_0>w_0/w_1\), accepting where it is smaller, and using any
constant randomization on equality. This proves simple-law optimality.
For this test both neighborhood risks attain their worst values at
the simple pair. Every competitor's worst weighted risk is at least
its risk at that pair. These two observations prove the robust
weighted minimax assertion. A zero weight gives the stated constant
decision directly.

For the maximum-loss problem consider the finite measure
\(M=L_0Q_0^{\otimes n}+L_1Q_1^{\otimes n}\).
Using the score transformed to \([0,1]\), choose an upper quantile
threshold and a constant boundary randomization whose rejection
integral under \(M\) is \(L_1\).
Such a threshold exists: for any probability CDF, the upper-tail
mass brackets every prescribed value between its strict and weak
tails at an appropriate quantile, and randomizing on the intervening
atom fills their difference. This follows directly from the inverse-CDF
construction in the preceding proof, or from the same continuity
argument as \cref{thm:p15-likelihood-threshold}.
The prescribed proportion is \(L_1/(L_0+L_1)\in(0,1)\).
The resulting identity
\[
 L_0E_{Q_0^{\otimes n}}\varphi+
 L_1E_{Q_1^{\otimes n}}\varphi=L_1
\]
is the required equality of error losses.

This score-threshold test is most powerful at its own simple null
size. A threshold at infinite likelihood can occur here only in
the perfect-discrimination case: the \(M\) mass there is
\(L_1Q_1^{\otimes n}(g_0=0)\leq L_1\), so filling \(L_1\)
requires that probability to be one and full rejection there.
It then has both errors zero. Other thresholds fall under the
ordinary Neyman--Pearson comparison, including its level endpoints.
Let their equal loss be \(r\). If a competitor had both simple losses
strictly less than \(r\), it would have smaller null rejection
probability and larger alternative power, contradicting that comparison.
Thus every test has maximum simple loss at least \(r\).
The least-favorable order attains both robust losses of the chosen
test at the pair, whereas a competitor's robust maximum dominates
its simple maximum. This proves the robust minimax result.
\end{proof}
\end{annexitem}

\begin{annexitem}{annex:p15.robust.sensitivity}{thm:p15-robust-sensitivity}{Pointwise contamination extrema and the clipping limit}
\begin{proof}
Linearity gives \(E_{aP+\varepsilon H}\varphi=aE_P\varphi+
\varepsilon E_H\varphi\). Every \(H\) expectation lies between the
pointwise infimum and supremum of \(\varphi\).
Conversely a Dirac probability at any \(x\in\mathsf X\) realizes
\(\varphi(x)\). Choosing values arbitrarily near a nonattained
supremum or infimum proves both formulas.
No nominal almost-everywhere convention can remove those points
from the contamination class.

If \(0<e\leq e'<1\), then
\[
 (1-e)P+eH=(1-e')P+
 e'\{(e/e')H+(1-e/e')P\}.
\]
The expression in braces is a probability, proving nesting.
For \(e=0\), \(P\) belongs to every larger class by choosing \(H=P\).
Variation balls are nested by their defining radius inequalities.
When the null class enlarges, the feasible set of level-\(\alpha\)
tests shrinks. When the alternative class enlarges, the infimum power
of each test can only decrease. Taking the maximum over feasible
tests proves the monotonicity of optimal worst-case power.

The normalization targets \(1/(1-\varepsilon_i)\) increase with their
fractions; monotonicity of \(F,G\) proves that \(u=c'\) increases with
\(\varepsilon_1\), and \(v=1/c''\) increases with \(\varepsilon_0\).
At zero fractions the specified choices \(u=v=0\) preserve these
inequalities, although uniqueness was not claimed there.
For equal \(e\uparrow e_*=d/(1+d)\), the targets tend to \(1+d\).
Both \(F(1)\) and \(G(1)\) equal \(1+d\). Since \(P_0\ne P_1\),
\(d>0\), so this target exceeds one and has its unique solution at one.
Continuity and monotonicity of the normalization functions therefore
give \(c'\uparrow1\) and \(1/c''\uparrow1\).

For a fixed ratio \(r<1\), eventually \(r<c'\), so the normalized
log score in the body is exactly zero. For \(r>1\), eventually
\(r>c''\), so it is exactly one. This also covers ratios zero and
infinity. Sum over a fixed finite sample to obtain the claimed
pointwise limit. If \(P_i(r=1)=0\), a finite union of null events
shows that such a sample avoids the excluded equality set almost
surely. Arbitrary contaminants need not avoid it, so no additional
uniform limiting power assertion follows.
\end{proof}
\end{annexitem}

\begin{annexitem}{annex:p15.robust.testing.calculation}{thm:p15-robust-testing-calculation}{All normalization masses, randomizations and risks}
\begin{proof}
The nominal ratios at \(a,b,c\) are \(1/6,1,6\).
With \(a_0=a_1=9/10\), the lower normalization equation, for
\(1/6<c'<1\), is
\[
 \frac9{10}\left(\frac9{10}+\frac35c'\right)=1.
\]
It gives \(c'=19/54\), which indeed lies in that interval.
Symmetry gives \(c''=54/19\).
The null base masses are \((54,27,9)/100\), and the remaining
contamination mass \(10/100\) is placed at \(c\).
The alternative base masses are \((9,27,54)/100\), and its remaining
mass is placed at \(a\). This proves the stated laws and their
membership in the full contamination neighborhoods.

At one-observation level \(1/10\), the largest score is at \(c\).
Randomization there with \(10/19\) spends
\((19/100)(10/19)=1/10\) null probability and yields power
\((54/100)(10/19)=27/95\).
For the nominal nonrandomized test, the contamination formula gives
\((9/10)(1/10)+(1/10)1=19/100\), proving failure of the unchanged
nominal size guarantee.

For two observations, multiply the coordinate scores.
Under \(Q_0\), the bottom, next, middle, next and top masses are
\[
 .54^2,\quad2(.54)(.27),\quad2(.54)(.19)+.27^2,\quad
 2(.27)(.19),\quad .19^2 .
\]
These equal the first row of the table divided by 10000.
The reversed coordinate probabilities give its second row.
Full rejection at the top spends \(361/10000\), leaving
\(639/10000\) of the target size. The next mass is \(1026/10000\),
so the boundary randomization is \(639/1026=71/114\).
The alternative power is
\[
 \frac{2916}{10000}\left(1+\frac{71}{114}\right)
 =\frac{8991}{19000}.
\]
All worst-case assertions follow from the already proved product
order and minimax theorem, not from inspecting only the displayed
nominal laws.

For general \(0<e<1/3\), solve
\((1-e)(.9+.6c')=1\), giving
\(c'=(1+9e)/\{6(1-e)\}\).
It lies in \((1/6,1)\) precisely on the stated range.
The base and excess masses then give the displayed \(Q_i\).
Here \(d=\tfrac12\sum_x|P_0(x)-P_1(x)|=1/2\);
the critical equal contamination is \(d/(1+d)=1/3\).

For the variation construction with \(\varepsilon=1/5\), the lower
transfer is \(\delta=1/10\). Increase \(q_1(a)\) from \(1/10\)
to \(1/5\), decreasing \(q_0(a)\) from \(3/5\) to \(1/2\);
make the opposite transfer at \(c\), and leave \(b\) unchanged.
The ratios become \(2/5,1,5/2\), so these modifications lie in exactly
the claimed clipped regions and solve both transfer equations.
Each signed difference has positive and negative masses \(1/10\),
so its variation norm is \(1/5\).
Randomizing with probability \(1/2\) at \(c\) gives null size
\((1/5)(1/2)=1/10\) and alternative power \((1/2)(1/2)=1/4\).

For the normal example, division of the two normalized densities gives
\(r(x)=e^{2x}\). Choose \(c'=e^{-2},c''=e^2\).
The integral of the displayed unnormalized null density is
\[
 P_0(X<1)+e^{-2}P_1(X\geq1)=\Phi(2)+e^{-2}/2=D.
\]
It is strictly larger than one: on \(x>1\),
\(e^{-2}p_1(x)>p_0(x)\), and this region has positive measure.
It is strictly smaller than \(\int\max(p_0,p_1)=2\Phi(1)\),
because its reciprocal clipping parameter \(e^{-2}\) is smaller
than one in the strictly increasing normalization function.
Thus \(0<\varepsilon_{\rm N}=1-1/D<1-1/\{2\Phi(1)\}\), the separation threshold
for these normal laws. Reflection supplies the alternative
normalization and the same contamination fraction.
The density formula is therefore exactly the contamination construction.
Its score exceeds one precisely for \(x>0\), and reflection makes the
two errors equal. Integrating the null density on that half-line gives
\[
 D^{-1}\{\Phi(2)-\Phi(1)+e^{-2}/2\}=1-\Phi(1)/D.
\]
The equalizer theorem proves its minimax property. Direct integration
under \(P_0\) gives the smaller nominal error \(1-\Phi(1)\).
\end{proof}
\end{annexitem}

\begin{annexitem}{annex:p15.robust.sampling.boundary}{thm:p15-robust-sampling-boundary}{Different joint laws with different guarantees}
\begin{proof}
Under independent per-coordinate contamination, the two observations
are independent Bernoulli\((e)\), so their probability of differing is
\(e(1-e)+(1-e)e=2e(1-e)\).
Under the displayed whole-sample contamination, the vector is either
\((0,0)\) or \((1,1)\), so that probability is zero.
The mechanisms therefore give different joint laws even with these
specified nominal and contaminating atoms.
They also define different classes: the independently contaminated law
assigns total mass \(1-(1-e)^2>e\) outside the nominal point \((0,0)\),
so it is not in the whole-sample class with radius \(e\).
Conversely the displayed whole-sample law is not any independent
product of identical coordinate laws when \(0<e<1\).

For the repeated-draw counterexample, draw \(Z\sim Q_0\) once and set
\((X_1,X_2)=(Z,Z)\). Both marginals lie in the null contamination class.
The displayed two-observation test rejects at \((c,c)\), accepts at
\((a,a)\) and \((b,b)\), and never encounters its randomized
score-\(54/19\) boundary. Its error is therefore
\(\Pr\{Z=c\}=19/100\).
The joint independence premise is exactly what the product guarantee
requires and this example violates.
\end{proof}
\end{annexitem}

\begin{annexitem}{annex:p15.robust.mlr.extension}{thm:p15-robust-mlr-extension}{The ordered-family bridge with varying supports}
\begin{proof}
First prove stochastic monotonicity from the stated density premise.
Let \(E\) be a Borel lower set: \(y\in E,x<y\) implies \(x\in E\).
For \(x\in E,y\notin E\) one has \(x<y\).
Integrate the assumed cross-product inequality over these pairs.
Tonelli applies to the nonnegative density products and gives
\[
 P_{\theta'}(E^c)P_\theta(E)
 \geq P_{\theta'}(E)P_\theta(E^c).
\]
Since the complementary probabilities sum to one, this is exactly
\(P_\theta(E)\geq P_{\theta'}(E)\).
The argument includes atoms and differing supports.

Let \(S=\{p_{\theta_0}+p_{\theta_1}>0\}\), with endpoint ratios assigned
zero and infinity on the respective one-sided supports.
The same pointwise cross-product condition implies that this
extended ratio is nondecreasing on \(S\).
Define on the entire line
\[
 r_*(x)=\sup\bigl(\{r(y):y\in S,\ y\leq x\}\cup\{0\}\bigr).
\]
It agrees with \(r\) on \(S\) and is nondecreasing.
Every monotone extended-real function is Borel measurable:
its strict superlevel sets are upper intervals.
Thus this is the required measurable monotone extension, including
gaps in the endpoint supports. Clip \(r_*\) and apply the source's
factor \(b\), when present, to obtain a globally nondecreasing score
with the same likelihood ratio under the endpoint pair.

For an interior score threshold, \(E=\{\ell<t\}\) is a lower set.
For \(\theta\leq\theta_0\), a contaminated null law obeys
\[
 R_\theta(E)\geq a_0P_\theta(E)
 \geq a_0P_{\theta_0}(E)=Q_0(E).
\]
For \(\theta\geq\theta_1\), the alternative obeys
\[
 R_\theta(E)\leq a_1P_\theta(E)+\varepsilon_1
 \leq a_1P_{\theta_1}(E)+\varepsilon_1=Q_1(E).
\]
The outer score thresholds give empty or full events as before.
In the variation case the same argument replaces these affine
bounds by \(P_\theta(E)-\delta\) and \(P_\theta(E)+\delta\).
Hence the complete stochastic inequalities hold for the two union
classes. The endpoint pair remains in those classes.
The product-order proof and all three minimax comparisons consequently
apply unchanged, with their now verified ordered-family premises.
For zero-radius endpoint neighborhoods, the order additionally implies
\(R_0(\ell=\infty)=0\) and \(R_1(\ell=0)=0\) throughout the union
classes. A tuple requiring the product \(0\cdot\infty\) is therefore
null under every allowed law, and a fixed assignment there preserves
all risks. The subsequent sequential proofs use only this score order,
pair membership and their stated moment premises, so apply to the
same union classes as asserted.
\end{proof}
\end{annexitem}

\begin{annexitem}{annex:p15.robust.sequential.errors}{thm:p15-robust-sequential-errors}{Affinity, ordered stopped paths and finite expected termination}
\begin{proof}
By Cauchy--Schwarz,
\(\eta=\int\sqrt{q_0q_1}d\mu\leq1\).
Equality would require the two square-root densities to be proportional
almost everywhere; their unit integrals then make them equal.
Since \(Q_0\ne Q_1\), \(\eta<1\).
The likelihood identities, including one-sided support, give
\[
 E_{Q_0}\sqrt\ell=\eta,\qquad E_{Q_1}\ell^{-1/2}=\eta.
\]
Here \(\sqrt0=0\), \(\infty^{-1/2}=0\); the opposite infinities occur
only on null sets of the corresponding law.
Apply the stochastic order to these monotone functions, first truncating
them and then using monotone convergence, to obtain
\[
 E_{R_0}e^{\gamma/2}\leq\eta,\qquad
 E_{R_1}e^{-\gamma/2}\leq\eta.
\]
In particular a null log score \(\gamma\) cannot be \(+\infty\) with positive
probability, nor an alternative log score \(-\infty\).
Under either hypothesis all partial sums are well defined almost
surely even if they are continued after a one-sided infinite increment.
The implemented test stops at that increment.
For \(S_n=\sum_{j=1}^n\gamma(X_j)\), independence factors these
nonnegative moments and yields
\(E_{R_0}e^{S_n/2}\leq\eta^n\) and
\(E_{R_1}e^{-S_n/2}\leq\eta^n\).
If \(N>n\), then \(A<S_n<B\).
Markov's inequality gives the two asserted bounds for \(\Pr(N>n)\).
If \(\eta=0\), both bounds show \(N=1\) almost surely.
For any nonnegative integer-valued possibly infinite \(N\),
\(N=\sum_{n\geq0}{\bf1}_{\{N>n\}}\); Tonelli gives
\(EN=\sum_{n\geq0}\Pr(N>n)\).
The displayed geometric tails make this finite under every allowed law,
and in particular prove almost-sure termination.

Couple the null score increments with those from \(Q_0\) by the
independent-uniform construction, and apply the increasing logarithm.
The contaminated null partial sums are no larger at every time.
If their path exits at the upper boundary first, the larger path
must already have exited at the upper boundary, or must do so by
that time. It cannot have exited at the lower boundary earlier:
at such a time the smaller path would also have crossed the lower
boundary, contradicting its first upper exit.
This proves inclusion of the null error event and therefore its
probability comparison.
For an alternative the order reverses; a first lower exit by the
larger path forces a lower exit by the smaller \(Q_1\) path no later.
This gives the type-II comparison. Membership of \(Q_i\) makes both
worst cases attained.
The comparison concerns exit decisions, not an ordering of the two
stopping times: a larger path can leave through either endpoint.
\end{proof}
\end{annexitem}

\begin{annexitem}{annex:p15.robust.sequential.asymptotic}{thm:p15-robust-sequential-asymptotic}{Uniform overshoot control and the stopped-sum identity}
\begin{proof}
Let \(C=\max(|\log\ell_-|,|\log\ell_+|)<\infty\).
Bounded positive scores make \(Q_0,Q_1\) equivalent, with
\(E_{Q_0}\ell=1\). The functions \(z-1-\log z\) and
\(z\log z-z+1\) are nonnegative, with equality only at one,
as follows by differentiation. Since the distinct pair has a
nonconstant ratio, integration gives
\[
 E_{Q_0}\gamma<0,\qquad
 E_{Q_1}\gamma=E_{Q_0}(\ell\log\ell)>0.
\]
The score order therefore supplies
\(E_{R_0}\gamma\leq-\kappa_0<0\) and
\(E_{R_1}\gamma\geq\kappa_1>0\), where
\(\kappa_0=-E_{Q_0}\gamma,\ \kappa_1=E_{Q_1}\gamma\).
For a null law \(R_0\), the stochastic order gives
\(E_{R_0}e^\gamma\leq E_{Q_0}e^\gamma=1\).
Independence makes \(e^{S_n}\), \(S_n=\sum_{j\leq n}\gamma(X_j)\),
a nonnegative supermartingale. To avoid assuming an optional-stopping
formula, condition one step at a time: on \(\{N\leq n\}\) the
stopped exponential stays fixed; on \(\{N>n\}\), independence and
the moment bound make its next conditional expectation at most
its current value. Induction yields
\(E_{R_0}e^{S_{N\wedge n}}\leq1\).
The already proved finite termination and Fatou's lemma imply
\(E_{R_0}e^{S_N}\leq1\). Therefore the upper-exit probability \(p\)
satisfies \(p\leq e^{-B}\).
For an alternative law, use
\(E_{R_1}e^{-\gamma}\leq E_{Q_1}e^{-\gamma}=1\) in the same proof.
Its lower-exit probability is at most \(e^A\).
These bounds are uniform over the neighborhoods.

For any fixed allowed law with mean \(m=E_R\gamma\), write
\[
 S_N=\sum_{j\geq1}\gamma(X_j){\bf1}_{\{N\geq j\}}.
\]
The absolute expected sum is at most
\(C\sum_j\Pr(N\geq j)=CEN<\infty\), by the preceding theorem.
Thus signed sum/integral interchange is justified.
The event \(\{N\geq j\}\) depends only on observations before \(j\),
so independence gives
\[
 E_RS_N=\sum_{j\geq1}m\,\Pr(N\geq j)=mE_RN.
\]
This proves the required stopped-sum identity with its actual
integrability premises.

At a lower exit \(S_N\in[A-C,A]\), and at an upper exit
\(S_N\in[B,B+C]\). Under a null law with upper-exit probability \(p\),
therefore
\(E S_N=A(1-p)+Bp+\zeta\), with \(|\zeta|\leq C\).
Since \(m<0\),
\[
 \left|\frac{E_RN}{|A|/|m|}-1\right|
 \leq p(1+B/|A|)+C/|A|
 \leq e^{-B}(1+B/|A|)+C/|A|\longrightarrow0.
\]
This bound is uniform in \(R_0\).
For the alternative the same calculation, with lower-exit probability
at most \(e^A\), gives the uniform relative-error bound
\[
 e^A(1+|A|/B)+C/B\longrightarrow0.
\]
Neither bound requires a relative rate of divergence of the two
boundaries.

Finally \(|E_{R_0}\gamma|\geq|E_{Q_0}\gamma|=\kappa_0>0\).
The uniform relative approximations bound
\(\sup_{R_0}E_{R_0}N\) above by
\((|A|/\kappa_0)(1+o(1))\), whereas \(E_{Q_0}N\) equals
\((|A|/\kappa_0)(1+o(1))\). Since \(Q_0\) is allowed, their ratio
is at least one and tends to one. Apply the identical reasoning to
\(E_{R_1}\gamma\geq E_{Q_1}\gamma=\kappa_1>0\) for the other ratio.
\end{proof}
\end{annexitem}
